\section{\textcolor{red}{Real-data examples: Gender gaps in ICU and Census data}}
%REV \section{Real-data example: Gender gap in ICU data}
\label{sec:real_example}

% Extended to the two Census outcomes (new cells in red). Original ICU-only table kept below as %REV.
\begin{table*}[t]
    \caption{Number of data subsets exhibiting at least one sign flip (in either the explained or unexplained component), by model\rc{, for the ICU data (out of $139$ subsets) and for the two outcomes of the U.S.\ Census data (out of $217$ and $183$ subgroups; \cref{appendix:us_labor_force})}. Columns $10\%$, $5\%$, and $1\%$ report the number of flipped subsets that are significant at the corresponding level. \rc{Dashes: logistic regression does not apply to the continuous log-income outcome.}}
    \label{tab:flip_summary_main}
    \centering
    \small
    \begin{tabular}{l cccc cccc cccc}
    \toprule
     & \multicolumn{4}{c}{ICU mortality} & \multicolumn{4}{c}{\rc{Census, log income}} & \multicolumn{4}{c}{\rc{Census, health insurance}} \\
    \cmidrule(lr){2-5} \cmidrule(lr){6-9} \cmidrule(lr){10-13}
    Model & Total flips & $10\%$ & $5\%$ & $1\%$ & \rc{Total flips} & \rc{$10\%$} & \rc{$5\%$} & \rc{$1\%$} & \rc{Total flips} & \rc{$10\%$} & \rc{$5\%$} & \rc{$1\%$} \\
    \midrule
    Linear      & 27 &  2 &  2 &  2 & \rc{33} & \rc{16} & \rc{13} & \rc{8}  & \rc{65} & \rc{38} & \rc{32} & \rc{23} \\
    Logistic    & 29 &  2 &  2 &  1 & \rc{--} & \rc{--} & \rc{--} & \rc{--} & \rc{62} & \rc{31} & \rc{25} & \rc{14} \\
    Neural net  & 62 & 29 & 17 &  6 & \rc{87} & \rc{--} & \rc{--} & \rc{--} & \rc{65} & \rc{--} & \rc{--} & \rc{--} \\
    XGBoost     & 62 & 17 & 10 &  3 & \rc{34} & \rc{18} & \rc{16} & \rc{12} & \rc{54} & \rc{34} & \rc{26} & \rc{16} \\
    TabPFN      & 76 & 22 & 19 &  5 & \rc{--} & \rc{--} & \rc{--} & \rc{--} & \rc{--} & \rc{--} & \rc{--} & \rc{--} \\
    \bottomrule
    \end{tabular}
\end{table*}

The index number problem reflects that the two reference choices for the (N)OBD could yield two different sets of values for the explained and unexplained components. We are not aware of an illustration in the literature that the two directions (for either the OBD or NOBD) can lead to \emph{substantively} different conclusions. We supply such \textcolor{red}{illustrations} next for \textcolor{red}{real-data analyses in healthcare and in the U.S.\ labor force}. %REV We supply such an illustration next for a real-data analysis in healthcare.

\textbf{Data.} \textcolor{red}{We consider two sources of data. The first is} the PhysioNet cohort %REV To construct an example, we examine data from the PhysioNet cohort
 of intensive care unit (ICU) patients collected for a study of in-hospital mortality \citep{goldberger2000physiobank}, where the outcome $Y$ is binary, indicating whether the patient died in the hospital or not. The dataset contains routine measurements taken at admission, including heart rate (HR), mean arterial pressure, temperature, and urine output. These values are recorded for every patient at the time of entry into the ICU and provide a description of their initial physiological state. To simulate various real-data analyses where practitioners might have access to only a subset of these patients, we construct $139$ versions of this dataset with various patient subsets; see \cref{appendix:icu_details}.
For reference groups, we consider males and females.
\textcolor{red}{The second is U.S.\ Census data: following \citet{bach2024heterogeneity}, we take the 2016 American Community Survey and consider two outcomes, the natural logarithm of income and a binary indicator for having health insurance, across $217$ and $183$ subgroups of workers defined by state and industry; see \cref{appendix:us_labor_force}.}

%Medical practitioners often look at moderate elevations in heart rate as a first indication of stress, infection, or concerning changes in blood flow and circulation. In the present example, we restrict our analysis to patients whose admission heart rate lies between the $50$th and $75$th percentiles. This patient set, which we refer to as HR quartile two, reflects a clinical presentation in which the heart rate is elevated but not extreme.

\textbf{Analysis.}
As discussed in \cref{sec:ob_real}, practitioners often draw conclusions based on the sign and statistical significance of the (N)OBD explained and unexplained components.
To explore whether these conclusions can be robust to the choice of reference group, we compute the explained and unexplained components for each \textcolor{red}{data subset} across one foundation model %REV for each of our $139$ data subsets
 and four classic models. The foundation model is TabPFN, a $25$ million parameter, transformer-based model, pre-trained to map training datasets to predictive distributions $Y|X$ \citep{hollmann2023tabpfn}. The other four models are linear regression,\footnote{Despite the binary outcome $Y$, it is common practice in applied health and policy research to fit separate linear models (via ordinary least squares) for the relationship between a binary outcome and covariates in each group, and then apply the OBD \citep{jann2008, edokaChangesCatastrophicHealth2017, sujin_LPM_OB, mweembaGapSelfRatedHealth2023}. In fact, this use of the linear model for a binary outcome is so common that it has a name, the ``linear probability model.''} logistic regression, XGBoost \citep{chen2016xgboost}, and a two-layer fully-connected neural network.
% \textcolor{red}{All models are trained and evaluated in-sample. We detail hyperparameter choices in \cref{appendix:icu_details}.}
%To examine a more modern, transformer-based model, we also consider TabPFN, a $25$ million parameter, pre-trained model trained to map training datasets to predictive distributions $Y|X$ \citep{hollmann2023tabpfn}.
We assess statistical significance using $1{,}000$ bootstrap replications.
\cref{tab:flip_summary_main} counts the number of sign flips across these models \textcolor{red}{and datasets} at various levels of statistical significance.
\textcolor{red}{In every dataset and for every model, we observe sign flips.}
\textcolor{red}{Within each dataset, the data subsets overlap and are not independent analyses. Our aim is an existence proof that sign flips arise in plausible analyses of real data, not an estimate of how frequently they arise in practice; for this reason, we do not apply a multiple testing correction to these counts.}
Concerningly, more complex models tend to have even more flips; for TabPFN \textcolor{red}{on the ICU data}, we see that more than half of the \textcolor{red}{subsets} show opposite signs %REV Concerningly, we see that \emph{every} model contains at least one sign flip, and more complex models tend to have even more flips; for TabPFN, we see that more than  half of the datasets show opposite signs
 for at least one of the explained or unexplained components, and $13.7\%$ show significant sign flips at the $5\%$ level.

\textcolor{red}{In the remainder of this section, we focus on the ICU data, which lets us examine a single analysis in detail; the Census results are in \cref{appendix:us_labor_force}.}

% Within the data set described above, we will compare two patient groups across hospitals (rather than hospital-specific patient groups as in our introductory example). In particular, we let the groups be $g \in \{\text{men}, \text{women}\}$. 
% The mortality rate for men exceeds the mortality rate for women by about $3.5$ percentage points; see \cref{tab:hr_q2_decomposition}.
% We would like to understand to what extent the observed mortality difference between men and women in HR quartile two might be attributable to observed characteristics upon admission (covariates) vs.\ medical care (outcome given covariates). 
% To that end, let $Y$ be a binary indicator for in-hospital mortality, and let $X$ be the vector of admission covariates comprising heart rate, mean arterial pressure, temperature, urine output, age, and ICU type.
% Despite the binary outcome $Y$, it is common practice in applied health and policy research to fit separate linear models (via OLS) for the relationship between a binary outcome and covariates in each group, and then apply the OBD \citep{jann2008, edokaChangesCatastrophicHealth2017, sujin_LPM_OB, mweembaGapSelfRatedHealth2023}.\footnote{In fact, this use of the linear model for a binary outcome is so common that it has a name, the ``linear probability model.''} We provide a detailed description of our data analysis --- and our search procedure for a data analysis exhibiting a substantive conclusion change --- in \cref{appendix:icu_details}. 

%REV \begin{table}[!ht]
%REV     \caption{Number of data subsets (out of $139$) exhibiting at least one sign flip (in either the explained or unexplained component), by model. Columns $10\%$, $5\%$, and $1\%$ report the number of flipped subsets that are significant at the corresponding level.}
%REV     \label{tab:flip_summary_main}
%REV     \centering
%REV     \small
%REV     \begin{tabular}{lcccc}
%REV     \toprule
%REV     Model & Total flips & $10\%$ & $5\%$ & $1\%$ \\
%REV     \midrule
%REV     Linear      & 27 &  2 &  2 &  2 \\
%REV     Logistic    & 29 &  2 &  2 &  1 \\
%REV     Neural net  & 62 & 29 & 17 &  6 \\
%REV     XGBoost     & 62 & 17 & 10 & 3 \\
%REV     TabPFN     & 76 & 22 & 19 & 5 \\
%REV     \bottomrule
%REV     \end{tabular}
%REV \end{table}

\subsection{A deeper dive on one data subset}

\cref{tab:flip_summary_main} shows that sign flips are common.
But what do these flips really mean in practice?
And are they merely an artifact of noise or instead a structural phenomenon?
As is often the case in studying complex machine learning models, studying linear regression provides a tractable proxy problem.
In the present study, linear regression is also of direct interest since the OBD (with linear regression) is widely used in practice.
So, for the rest of the paper, we focus on \cref{eq:OB_H_Counterfactual,eq:OB_K_Counterfactual} with linear regression (i.e., the OBD instead of NOBD).
To start answering our questions, we examine one particular data subset leading to a significant sign flip for linear regression in \cref{tab:flip_summary_main}.
We will see what it means for sign flips to change meaningful conclusions and also see that flips have a natural structural interpretation.

The particular data subset we consider covers patients with an elevated heart rate.
Medical practitioners often look at moderate elevations in heart rate as a first indication of stress, infection, or concerning changes in blood flow and circulation. In the present example, we consider patients whose admission heart rate lies between the $50$th and $75$th percentiles. This patient set, which we refer to as HR quartile two, reflects a clinical presentation in which the heart rate is elevated but not extreme.

\textbf{Two different conclusions.} 
We report the OBD for each reference group in \cref{tab:hr_q2_decomposition}. We next interpret the substantive conclusions
we might draw from each reference group. First, consider women as the reference. The explained component is positive ($0.021$), with a statistically significant difference from zero.
This result suggests that differences in admission covariates place men at higher mortality risk. A medical practitioner might therefore examine whether various factors before admission differ across gender; these factors include early differences in blood flow and circulation, underlying disease burden, or overall severity at admission.

Second, consider men as the reference. The explained component is negative ($-0.007$), though the difference from zero is not statistically significant. In this case, the observed covariate differences seem to predict \emph{lower} (or equal) mortality for men relative to women rather than higher. In other words, the same covariates would now be interpreted as offering men a slight mortality advantage. This interpretation is substantively different than the interpretation that results from using women as the reference.

%\begin{table}[h]
%\caption{Oaxaca--Blinder decompositions for HR quartile two. We compute standard errors (in parentheses) via bootstrap resampling with $1{,}000$ replications. 
%    We calculate p-values using a normal approximation (Wald test), where the test statistic is the coefficient divided by its bootstrap standard error. 
%    Significance levels: * $p<0.10$, ** $p<0.05$, *** $p<0.01$.}
%\label{tab:hr_q2_decomposition}
%    \begin{center}
%    \begin{small}
%    \begin{sc}
%    \begin{tabular}{lccc}
%    \toprule
%    \textbf{Reference} & \textbf{Explained} & \textbf{Unexplained} & \textbf{Total gap} \\
%    \midrule
%    %\hline \\[-6pt]
%    Women 
%    & \makecell{$0.021^{***}$ \\ $(0.007)$}
%    & \makecell{$0.014$ \\ $(0.025)$}
%    & \makecell{$0.035$ \\ $(0.023)$}
%    \\
%    Men 
%    & \makecell{$-0.007$ \\ $(0.011)$}
%    & \makecell{$0.042^{*}$ \\ $(0.025)$}
%    & \makecell{$0.035$ \\ $(0.023)$}
%    \bottomrule
%    \end{tabular}
%    \end{sc}
%    \end{small}
%    \end{center}
%\end{table}

\begin{table}[!ht]
  \caption{Oaxaca--Blinder decompositions for HR quartile two. The total gap is men's mean minus women's mean. We compute standard errors (in parentheses) via bootstrap resampling with $1{,}000$ replications. 
    We calculate p-values using a normal approximation (Wald test), where the test statistic is the coefficient divided by its bootstrap standard error. 
    Significance levels: * $p<0.10$, ** $p<0.05$, *** $p<0.01$.}
  \label{tab:hr_q2_decomposition}
  \begin{center}
    \begin{small}
      %\begin{sc}
        \begin{tabular}{lcccr}
          \toprule
          Reference & Explained & Unexplained & Total gap  \\
          \midrule
              Women 
    & \makecell{$0.021^{***}$ \\ $(0.007)$}
    & \makecell{$0.014$ \\ $(0.025)$}
    & \makecell{$0.035$ \\ $(0.023)$}
    \\
    Men 
    & \makecell{$-0.007$ \\ $(0.011)$}
    & \makecell{$0.042^{*}$ \\ $(0.025)$}
    & \makecell{$0.035$ \\ $(0.023)$}
    \\
          \bottomrule
        \end{tabular}
      %\end{sc}
    \end{small}
  \end{center}
  \vskip -0.1in
\end{table}


\textbf{How does the sign flip arise?}
%Recall that the explained component is given by $\EHOB := \Delta\mu^T \beta_H$ when $H$ is the reference group, and $\EKOB := \Delta\mu^T \beta_K$ when $K$ is the reference group.
To understand how the sign flip arises, it will help to introduce more notation for the linear regression case.
We now restrict to real-valued covariates $X \in \mathbb{R}^d$, and we assume we fit a linear regression to each group $g$: $\alpha_g + X^\top \beta_g$
with intercept $\alpha_g \in \mathbb{R}$ and slope $\beta_g \in \mathbb{R}^d$. We let
$\mu_g := \E_g[X],
\Delta \mu := \mu_H - \mu_K,
\Delta \beta := \beta_H - \beta_K,
\Delta \alpha := \alpha_H - \alpha_K,
$
and use a hat to denote estimators of these quantities (cf.\ \cref{sec:ob_real}).

With these choices, the OBD
(\cref{eq:OB_H_Counterfactual,eq:OB_K_Counterfactual}) becomes
\begin{equation}
    \Delta_Y = \underbrace{\Delta \mu^\top \beta_H}_{E^{(H)}} + \underbrace{\mu_K^\top \Delta \beta + \Delta \alpha}_{U^{(H)}} = \underbrace{\Delta \mu^\top \beta_K}_{E^{(K)}} + \underbrace{\mu_H^\top \Delta \beta + \Delta \alpha}_{U^{(K)}}.
\end{equation}
The first equality is the OBD with reference group $H$, and the second with reference group $K$.

With this notation in hand, we observe that $\Delta\hat{\mu}$ is fixed between $\hat{E}^{(H)}$ and $\hat{E}^{(K)}$. So it must be that some components of $\hat{\beta}_H$ and $\hat{\beta}_K$ are changing sign to give $\hat{E}^{(H)} := \Delta\hat{\mu}^T \hat{\beta}_H$ and $\hat{E}^{(K)} := \Delta\hat{\mu}^T \hat{\beta}_K$ opposite signs.
In \cref{fig:sign_flip_overview} (left), we show the fitted linear models of in-hospital mortality on admission heart rate, holding all other covariates fixed at their group means.
We see that the signs of the slopes differ.
In Appendix \cref{tab:delta_mu_betas}, we report values for all elements of each vector $\hat{\beta}_{g}$ in the full analysis. 
We see that there are multiple covariates $c$ for which $\hat{\beta}_{\textrm{men},c} < 0 < \hat{\beta}_{\textrm{women},c}$.
And as $\Delta\hat{\mu}_c$ has the same sign for all of these covariates, each one pushes towards a sign flip. %, and additionally, $\Delta\mu$ is positive for all of these covariates.
In aggregate, these changes add enough to cause a sign flip in the explained component.
%In \cref{sec:simulated_example}, we show that these phenomenon are not artifacts of finite-samples or model misspecification; in fact, we show that sign flips can occur for well-specified linear models with infinite data.

\begin{figure}[!ht]
    \centering
    \begin{subfigure}[t]{0.48\linewidth}
        \centering
        \includegraphics[width=\linewidth]{Sections/Figures/hr_quartile2_mortality_lines.pdf}
        \label{fig:hr_quartile2}
    \end{subfigure}
    \hfill
    \begin{subfigure}[t]{0.48\linewidth}
        \centering
        \includegraphics[width=\linewidth]{Sections/Figures/prob_unexplained.pdf}
        \label{fig:prob_unexplained}
    \end{subfigure}
    \caption{Left: Linear fit of mortality on admission heart rate for each group when holding all other covariates fixed at their group means. Right: Percentage of parameter space leading to sign flips. See \cref{remark:computing_irwin_hall} for how we compute the percentage for the unexplained component.}
    \label{fig:sign_flip_overview}
\end{figure}


\textbf{Limitations.}
Our analysis reveals sign flips in the explained and unexplained components of the (N)OBD in a real dataset.
It is a priori possible that some of our choices may be unusual; these choices include the dataset itself, our choice of covariates, our choice of outcome, or our particular subset of data. And it is possible that sign flips may not occur in other settings.
To address this possibility, we provide a second illustration analyzing differences in income as well as health insurance status using U.S.\ Census data in \cref{appendix:us_labor_force}.
The latter data have previously been analyzed by \citet{bach2024heterogeneity}, reducing the concern that our empirical results are due to atypical choices on our part.
We also give a theoretical account of sign flips in \cref{sec:signflips}, where we establish that they correspond to a nontrivial region of the parameter space.